\chapter{Background}\label{ch:background}

This chapter establishes the background needed to understand the design of DUET. It begins with
Signal's asynchronous key-distribution model and then examines key-transparency systems and
continuous-authentication research to derive the requirements carried into \cref{ch:design}.
\Cref{app:extended-background} provides further detail on protocol mechanics, usability evidence, and
the systems reviewed in this chapter.

% ============================================================================
\section{The Signal Protocol Stack}\label{sec:bg-signal}

Signal is a widely used end-to-end encrypted messenger~\cite{signal-home,dowling2020nocompromise}. Its
protocol combines X3DH for asynchronous session establishment with the Double Ratchet for ongoing key
evolution~\cite{x3dh,doubleratchet}. The messaging server supplies the public pre-key material needed to start a
conversation while the recipient is offline~\cite{signaldocs}, so confidentiality depends on the
authenticity of keys obtained from provider-operated infrastructure.

\subsection{Key Material, Safety Numbers, and What They Cover}\label{ssec:bg-keymat}

Signal publishes a long-term identity key, a signed pre-key, and optional one-time pre-keys. The Safety
Number authenticates the parties' identity keys and account identifiers through an out-of-band
comparison~\cite{safetynumbers}. A signed pre-key is authenticated indirectly by its identity-key
signature, while one-time pre-keys are unsigned. The ceremony therefore authenticates the long-term
identity binding rather than every public value used in one session. \Cref{sec:ext-signal} provides the
full key-material notation and coverage.

\subsection{The Extended Triple Diffie--Hellman (X3DH) Protocol}\label{ssec:bg-x3dh}

X3DH lets Alice establish a shared secret $\sharedS$ with an offline Bob by combining identity,
ephemeral, signed-pre-key, and optional one-time-pre-key Diffie--Hellman terms~\cite[Section~2.4]{x3dh}.
The resulting secret initializes the later protocol stages and is also the input from which DUET
derives conversation labels. Alice verifies Bob's signed pre-key under the supplied identity key and
aborts if that signature fails~\cite[Section~3.3]{x3dh}.

When Bob's genuine one-time pre-key is absent or withheld, an attacker that substitutes Bob's identity
key and a signed pre-key valid under that identity can derive the resulting secret and impersonate
Bob. Replacing only the signed pre-key fails the signature check, while replacing only the one-time
pre-key normally causes the parties to derive different secrets. DUET targets the coordinated
substitution, whose full derivation and message sequence are in \cref{sec:ext-signal}.

\subsection{The Double Ratchet Protocol}\label{ssec:bg-ratchet}

The Double Ratchet evolves the X3DH or PQXDH secret into per-message keys, providing forward secrecy
and post-compromise recovery after the adversary becomes passive~\cite{doubleratchet,pqxdh}. It does
not re-authenticate the public keys used to establish the session, so its guarantees still depend on
an authentic initial binding. \Cref{sec:ext-ratchet} explains the key-schedule details.

\subsection{Safety Numbers and the Verification Gap}\label{ssec:bg-gap}

Signal provides manual Safety Number comparison and optional Automatic Key Verification. AKV checks a
peer's stored binding against Signal's transparency service, but its documented availability is
limited for some username and phone-number-discovery cases, and it does not provide a channel through
which a detecting client can warn the peer~\cite{signal-akv}.

Usability evidence also limits reliance on manual verification. Across unguided trials in one study,
only 14\% of pair--application trials succeeded, and only two of twelve pairs completed the ceremony
in any tested application~\cite[p.~30]{vaziripour2017}. Instructions improved
success~\cite[p.~29]{vaziripour2017}, but users still needed substantial
time~\cite[p.~36]{vaziripour2017}. ACKA and its predecessor MoDUSA authenticate evolving protocol
state, yet both retain a user-initiated comparison or an assumed secure
log~\cite{dowling2021meta,dowling2023acka}. The supporting measurements are reported in
\cref{sec:ext-usability}.

Three requirements follow. Authentication should detect coordinated substitution of
session-establishment keys, detection should not depend on a rarely completed ceremony, and a detected
substitution should reach the peer rather than remain local to the key owner.

% ============================================================================
\section{Secure Transparency Logs}\label{sec:bg-logs}

Key transparency makes the directory accountable by binding each response to a committed state. An
operator that substitutes a key must place it in a state whose later evolution can be monitored.

\subsection{Commitments and Merkle Log Trees}\label{ssec:bg-merkle}

A Merkle root commits to every leaf under a collision-resistant hash. An inclusion proof shows that a
record belongs to the committed tree, while a consistency proof shows that a later log is an
append-only extension of an earlier one~\cite{rfc6962,rfc9162}. Both proofs are logarithmic in the
relevant tree size, and they answer different questions, because inclusion authenticates one record
whereas consistency shows that committed history was not rewritten. The construction, domain
separation, and proof walkthrough are explained in \cref{sec:ext-merkle}.

Although each proof is compact, requesting one proof for every missed epoch makes a returning
client's total work grow with its offline period. DUET's windowed proof addresses this cost.

\subsection{Key Transparency Requirements}\label{ssec:bg-ct2kt}

Key transparency applies verifiable commitments to identity-to-key bindings~\cite{chase2016overlays,melara2015coniks}.
A verifiable random function maps an identity to a deterministic but outsider-unpredictable label and
proves that the mapping was evaluated correctly~\cite{rfc9381}. This resists outsider enumeration, but
it does not hide queries from the operator or conceal the directory's scale.

Local proof verification also does not establish that every observer received the same state. A malicious
operator can present Alice and Bob with different internally consistent histories, and exposing such a fork
requires signed commitments obtained from independent vantage points to be compared. DUET therefore
separates local lookup verification from cross-vantage equivocation detection. The split-view problem
and the comparison mechanisms used by existing systems are discussed further in
\cref{sec:ext-splitview}.

\subsection{The IETF KeyTrans Architecture}\label{ssec:bg-keytrans}

The IETF Key Transparency architecture combines a VRF-indexed prefix tree with a log of committed
directory states and defines search, monitoring, and auditing roles~\cite{ietf-keytrans-architecture,ietf-keytrans-protocol}.
Signal's service implements the emerging protocol with modifications~\cite{signal-kt-server}. The draft
calls for reasonable monitoring windows but does not prescribe their construction. DUET uses the same
combined-tree model and supplies a windowed consistency mechanism.

% ============================================================================
\section{Verifiable Key Directories for Key Transparency}\label{sec:bg-vkd}

A verifiable key directory is an authenticated key--value map whose membership and non-membership
answers can be checked against a public commitment. SEEMless formalized completeness, soundness, and
simulation-based privacy for this abstraction~\cite{chase2019seemless}, and later constructions refine
its update history, storage, monitoring, and privacy properties.

\subsection{Cryptographic Building Blocks}\label{ssec:bg-vkd-blocks}

Most systems combine VRF-derived labels with a sparse or path-compressed Merkle directory, and some add
a zero-knowledge set to hide information beyond the queried binding. SEEMless uses an append-only ZKS,
Parakeet an ordered ZKS, and ELEKTRA a rotatable ZKS~\cite{chase2019seemless,malvai2023parakeet,chen2022rzks,len2023elektra}.
Merkle\textsuperscript{2} instead composes a chronological log with per-key
histories~\cite[pp.~1--2]{hu2021merkle2}. \Cref{sec:ext-vkd-blocks} explains the structural
differences and terminology.

\subsection{Academic and Deployed Systems}\label{ssec:bg-vkd-academic}

The literature develops a sequence of trade-offs. CONIKS introduced VRF-indexed bindings and
epoch-by-epoch monitoring~\cite{melara2015coniks}. SEEMless formalized the directory abstraction,
Parakeet added storage compaction and witnesses, ELEKTRA addressed multi-device histories and privacy
recovery, OPTIKS reduced history cost, and PRISM moved towards whole-directory verification~\cite{chase2019seemless,malvai2023parakeet,len2023elektra,optiks2024,pusch2025prism}.
Detection in each still depends on a user, an auditor, a witness set, or a global verifier rather than
a conversation-specific warning delivered to the peer.

Two results sharpen that comparison. ELEKTRA's unsuppressible warning is raised after a full account
reset rather than after substitution of one key~\cite[p.~2]{len2023elektra}. PRISM frames the wider
problem as reliance on someone else to check, whether a user, gossiping clients, a blockchain, or an
auditor~\cite[p.~6]{pusch2025prism}. Its whole-directory proof changes who verifies the state but does
not itself provide conversation-specific peer notification.

CONIKS is the principal baseline for this thesis because it combines private VRF-derived labels with
signed epoch roots and self-monitoring. A client verifies its own binding and later exposes a fork by
comparing signed roots with another observer. Its limitations are also central here, because
monitoring work grows with the offline period, direct client gossip was left as future work, and
detection does not itself notify the communication peer~\cite{melara2015coniks}. SEEMless reduces
history cost through versioned labels, Parakeet replaces external gossip with a witness quorum, and
ELEKTRA adds signed device-key histories. Each improvement introduces a different monitoring, trust, or
deployment condition rather than removing the peer-warning gap.

Production systems demonstrate that these constructions can operate at messaging scale. WhatsApp,
Signal, Proton, Google, and Apple use different combinations of monitoring, auditing, and commitment
comparison~\cite{whatsapp2023whitepaper,signal-kt-server,goebel2023proton,google-kt,apple-ckv}, and an
independent review of the Meta implementation reports no break of its transparency
guarantees~\cite{nccgroup2023akd}. None of the cited deployments provides the peer-warning channel
defined in this thesis. The per-system mechanisms, limitations, and citations are examined further in
\cref{sec:ext-vkd-survey} and summarized in \cref{tab:bg-compare}.

The deployments also expose relevant engineering choices. WhatsApp demonstrates a public audit path at
messaging scale but disables client history monitoring in the cited configuration~\cite{whatsapp2023whitepaper}, Signal follows the
IETF architecture and uses fixed-length prefix copaths~\cite{signal-kt-server}, and Proton combines owner and
external auditing with Certificate Transparency. Proton also identifies communication of verification
failures to users as unresolved~\cite[p.~88]{goebel2023proton}. These systems establish deployment feasibility, but
their detection remains local to the account owner or an external verifier.

% ============================================================================
\section{The ACKA (Authenticated Continuous Key Agreement) Framework}\label{sec:bg-acka}

ACKA composes continuous key agreement, Continuous Entity Authentication (CEA), and an abstract secure
log~\cite[p.~18]{dowling2023acka}. CEA derives a label and a fingerprint from chained secret state and
the new public key-share at each epoch~\cite[pp.~13--14]{dowling2023acka}, and its security notion is
$\CEAsec$~\cite[p.~18]{dowling2023acka}. A forged share either collides at the honest label, producing
two fingerprints at one label, or appears under a label the impersonated party never derives and is
caught on lookup~\cite[pp.~7--8]{dowling2023acka}. The secure-log scheme $\Log$, required to satisfy
$\Logsec$, stores and returns those fingerprints~\cite[pp.~10--11]{dowling2023acka}.
\Cref{fig:bg-acka} summarizes the composition.

\begin{figure}[H]
\centering
% fig-bg-acka.tex — extracted from ch2-background.tex
% Label: \label{fig:bg-acka}
\begin{tikzpicture}[>=Stealth, font=\small,
   box/.style={draw, rounded corners, minimum width=26mm, minimum height=20mm,
               align=center, text width=24mm, inner sep=2pt},
   open/.style={draw, densely dashed, rounded corners, minimum width=26mm,
                minimum height=20mm, align=center, text width=24mm, inner sep=2pt}]
  \node[box]  (cka) at (0,0)    {\textbf{CKA}\\[2pt]\footnotesize continuous key agreement (the ratchet)};
  \node[box]  (cea) at (5.6,0)  {\textbf{CEA}\\[2pt]\footnotesize per-epoch label \& fingerprint (security: CEAsec)};
  \node[open] (log) at (11.2,0) {\textbf{LOG}\\[2pt]\footnotesize append-only, immutable log (security: Logsec)};
  \draw[->] (cka) -- node[above,font=\scriptsize]{public $T_t$}
                     node[below,font=\scriptsize]{secret $I_t$} (cea);
  \draw[->] ([yshift=2.5mm]cea.east) -- node[above,font=\scriptsize]{post} ([yshift=2.5mm]log.west);
  \draw[<-] ([yshift=-2.5mm]cea.east) -- node[below,font=\scriptsize]{query} ([yshift=-2.5mm]log.west);
  \node[align=center, text=black!55, font=\footnotesize] at (11.2,-1.9)
       {\emph{concrete realization left open}};
\end{tikzpicture}

\caption[ACKA composition]{ACKA composes CKA and CEA with the abstract
secure-log interface $\Log$, whose required security notion is $\Logsec$. Adapted from Dowling and
Hale~\cite{dowling2023acka}.}
\label{fig:bg-acka}
\end{figure}

ACKA addresses an active channel adversary and provides continuous authentication, but its detection
result is conditional on a log satisfying $\Logsec$ rather than on a directory that resists a
malicious operator~\cite[p.~7]{dowling2023acka}. ACKA assumes a common initial state associated with
trust on first use and verifies sequentially by epoch~\cite[p.~6]{dowling2023acka}. Only the receiver
learns of a forgery, and its consistency mechanism is not used for third-party fork detection. The
formal CKA definition, adversary model, abstract log construction, and receiver-only detection result
are examined further in \cref{sec:ext-acka}.

DUET supplies the verifiable directory behind $\Log$. It also adds conversation-derived lookup labels,
a shared peer-warning slot, windowed catch-up, and cross-vantage comparison of signed heads. ACKA and
key transparency therefore address complementary adversaries. ACKA authenticates evolving state
against an active channel attacker, while key transparency constrains a malicious directory operator.

% ============================================================================
\section{Comparative Analysis of Existing Systems}\label{sec:bg-compare}

Key-transparency systems differ in the objects they authenticate, the mechanisms they use to detect
misbehaviour, their support for offline recovery, and their consistency guarantees.
\Cref{tab:bg-compare} summarizes these differences.

\begin{table}[t]
\centering
\caption[Cross-system comparison of key-transparency directories.]{Cross-system comparison of
key-transparency directories across their authenticated objects, detection, offline recovery, and
consistency.}
\label{tab:bg-compare}
\footnotesize
\setlength{\tabcolsep}{3pt}
\renewcommand{\arraystretch}{1.0}
\begin{tabular}{@{}>{\raggedright\arraybackslash}p{2.0cm}>{\raggedright\arraybackslash}p{2.3cm}>{\raggedright\arraybackslash}p{2.5cm}>{\raggedright\arraybackslash}p{1.35cm}>{\raggedright\arraybackslash}p{2.15cm}>{\raggedright\arraybackslash}p{2.6cm}@{}}
\toprule
System & What it logs & Primary detection mechanism & In-band peer warning &
  Offline recovery & Consistency and fork resistance \\
\midrule
CONIKS~\cite{melara2015coniks} & username $\to$ key binding & self-monitoring & No & linear in the number of missed updates &
  STR gossip among auditors\newline direct client-to-client exchange left as future work \\
SEEMless~\cite{chase2019seemless} & identity key (versioned) & self-monitoring (client history checks) & No & per-update
  history & $\ge$1 honest auditor per epoch\newline commitments compared through gossip or a blockchain \\
Parakeet~\cite{malvai2023parakeet} & identity key & self-monitoring & No & bounded by the compaction interval & $2f{+}1$ signatures from $3f{+}1$
  witnesses \\
ELEKTRA \cite{len2023elektra} & multi-device keychains & self-monitoring & Warning only after account reset & per-update &
  external auditor \\
OPTIKS~\cite{optiks2024} & identity key (+ devices) & self-monitoring (single proof, common case) & No &
  verify before old periods are deleted & $\ge$1 auditor + bulletin board \\
PRISM~\cite{pusch2025prism} & account keys & whole-directory verification & No &
  constant-size proof & zk-rollup on a data-availability layer \\
WhatsApp (AKD)~\cite{whatsapp2023whitepaper} & identity key & public auditing & No & Not applicable in the cited configuration\newline self-monitoring disabled & signed root + public write-once (WORM) audit log \\
Signal KT~\cite{signal-kt-server,signal-akv} & account identifier $\to$ identity key, phone-number and username-hash mappings & self-monitoring and audit interfaces, optional user-triggered AKV & No &
  monitoring endpoint\newline offline-window policy not documented in the cited source & signed head and public audit endpoints\newline independent auditor deployment not documented in the cited source \\
Proton KT~\cite{goebel2023proton} & email address $\to$ Signed Key List (SKL) fingerprint & self-audit and external audit & No &
  self-audit within 90 days & CT logs and external auditing \\
Google KT~\cite{google-kt} & account keys & self-monitoring (experimental, archived 2024) & No & Not documented in the cited source & CONIKS map in a CT log \\
Apple CKV~\cite{apple-ckv} & account keys, device keys, and opt-in state & automatic device verification and E2E-encrypted CloudKit
  cross-check & No & Not documented in the cited source & on-device consistency + encrypted log-hash gossip + internal audit \\
\bottomrule
\end{tabular}

{\scriptsize\raggedright \emph{In-band peer warning} means a conversation-specific substitution
warning written by one party and verified by the other. SEEMless additionally assumes that a user
recalls their own update times. The in-band mechanisms in Apple CKV and PRISM carry commitments rather than peer-substitution
warnings. Merkle\textsuperscript{2} and ACKA are omitted because neither defines a complete
identity-to-key directory. ACKA instead assumes the secure-log interface described in
\cref{sec:bg-acka}.\par}
\end{table}

\FloatBarrier
\subsection{Implications for the Design}\label{ssec:bg-implications}

The comparison yields three requirements. First, detection must reach the peer. Existing systems assign
discovery to the key owner, an auditor, a witness set, or a global verifier, while the cited WhatsApp
configuration falls back to manual comparison after a failed
check~\cite[pp.~5, 10]{whatsapp2023whitepaper}. ELEKTRA warns only after a full account
reset~\cite[p.~2]{len2023elektra}. DUET
therefore needs a conversation-specific warning that remains available while either party is offline.

Second, catch-up must not require one proof per missed epoch. The IETF draft requests reasonable
monitoring windows without defining the mechanism~\cite{ietf-keytrans-protocol}. CONIKS identifies
linear offline monitoring as an open cost~\cite[p.~11]{melara2015coniks}.

Third, equivocation must be detectable outside the affected session. Locally valid proofs cannot show
whether another client received a different history, so signed heads must be compared across
independent vantage points. These requirements lead directly to the dual-index warning, windowed
consistency, and auditor mechanisms in \cref{ch:design}.

Proof construction introduces a separate bandwidth--privacy trade-off. Signal's fixed-length
\num{8192}-byte proof conceals directory occupancy~\cite{signal-kt-server}, whereas CONIKS transmits
the authentication path required for the queried entry~\cite[p.~10]{melara2015coniks}. The evaluation
places DUET's compressed proof within this design space.
